% ------------------------------------------------------------------------
% -*-TeX-*- -*-Hard-*- Smart Wrapping
% ------------------------------------------------------------------------


\chapter{結論}
\label{cha:6}

\section{まとめ}
\label{sec:6.1}

本研究は，4B級小型言語モデルにLoRAでペルソナを焼き込んだ3体のエージェント集団を構成し，議論によるチーム推論とその進化的最適化を，計算量を揃えたSelf-Consistencyという厳しい基準に対して検証した．得られた知見は以下の4点に要約される．

第一に，\textbf{議論の効果はベンチマーク依存であり，「多様性の利得」と「ペルソナ注入による能力毀損」の2軸のトレードオフとして整理できる}．素の議論は数学（MATH-500，$+4.2$pt）と高難度科学（SuperGPQA，$+3.1$pt）で有意に有効だが知識推論（MMLU-Pro，$-4.1$pt）で有意に有害であり，重みペルソナチームはその逆の符号を示した．この反転は「議論は多数決に勝てない」とする近年の理論的整理と「小型モデルでも議論が効く」とする報告の併存を説明しうる．ただし反転を「問題の検証可能性」に帰する解釈は仮説にとどまる——MMLU-Pro内部のカテゴリ別分析（\S\ref{sec:4.6.6}）では定量系カテゴリでも議論は有効化せず，反転の規定要因がベンチマーク構成（形式・難易度・選択肢の有無）の複合である可能性を排除できない．

第二に，\textbf{チーム誤差は個体能力損失と集約損失に分解でき，それぞれに実効的な処方が存在する}．ペルソナSFTは短い応答スタイルの刷り込みによって思考連鎖を半減させ長い導出能力を選択的に破壊していたが，ベースモデル自身が生成した検証済み長CoTのリプレイ混合（約38\%）を含む再学習によりベース水準（等価マージン$\pm 2$ptの同等性が成立）まで回復し，一部はベース単体を上回った（ただし再学習ではrank・学習率等も同時に変更しており，リプレイ単独の寄与は分離できていない）．集約損失（正解を持つ個体の答えをチームが採用できない損失，最大13pt）に対しては，logprob重み付き投票・条件付き更新・発言匿名化が小幅の改善を与えたが，大部分は未回収の課題として残る．

第三に，\textbf{チームレベル適応度による進化的最適化は，本実験の実装では機能する状態になかった}．進化前後の$\Delta W$はノルム比1.000かつcosine類似度0.9997--1.0000でほぼ同一であり，進化演算子は重み空間をほとんど動かしていなかった（探索の不全）．fitness sharingは$K=2$構成で選抜中立であり（多様性圧の不全），適応度測定の標準誤差（100問で$\pm 4.9$pt）に起因する勝者の呪いが選抜を無情報化していた（選抜の不全）．開発セットでの見かけの改善（$+12$pt）は最終評価に汎化せず，MATH-500で観測された$+4.3$ptも極小の重み摂動に対する生成過程の感度として説明するのが最も保守的である．したがって本研究は枠組み自体の有効性を支持も反証もしておらず，能力回復という目的にはリプレイ再学習が確実な処方であった．

第四に，診断に基づく処方（リプレイ再学習$\times$条件付き議論$\times$匿名化$\times$重み付き投票）を実測ゲート方式で積み上げた最終チームは，測定環境を統一した最終評価において\textbf{旧チームに$+2.7$pt（問題IDクラスタ解析，$p=0.0003$）と有意に優り，ペルソナSFTが生んだ能力毀損をベース単体との同等性（等価マージン$\pm 2$pt）が成立する水準（MATH-500では$+3.4$pt有意優位）まで回復させた}．しかし\textbf{生成回数を揃えたSC@9には$-3.5$pt（$p<10^{-4}$）と有意に及ばず}，「対策を尽くしても4B級のチーム議論は同等の生成予算の多数サンプリングに勝てない」という統制された負の結果が本研究の主結論である．付随して，評価コンテナイメージ間の系統差（同一設定で$+6$pt級，ベース素生成に選択的に作用しチーム側を見かけ上優位にするバイアス）を検出・統制した過程は，小型LLM評価の再現性に対する方法論的警鐘である．

\section{今後の課題}
\label{sec:6.2}

\begin{itemize}
\item \textbf{サンプリングと議論のハイブリッド}: 本研究のチームは1問あたり6生成とSC@9より少ない計算量で比較したが，各エージェントがまずSelf-Consistency@3で自己集約してから議論する構成（3体$\times$3サンプル$=$9生成）はSC@9と計算量が完全に一致し，独立サンプルの分散削減と議論の誤り訂正を同時に得られる可能性がある．本研究の実測（チーム化ゲインと集約損失の分解値）に基づけば，これが4B級でSC@9に到達しうる残された最有力の設計である．
\item \textbf{集約損失の本格的な回収}: oracle上限との差は依然大きい．検証器の学習（RLVRによる集約器訓練）や，正解が少数派の問題を対象とした再帰的自己集約の導入が有望である．
\item \textbf{選抜として機能する進化の再検証}: 識別力の高い問題の厳選（IRTベースの縮約評価）と逐次半減法により選抜ノイズを抑えた上で，チームレベル適応度が誤り相関の低い個体を選抜できるかを検証すること．本研究で保存した設計（研究設計書v3 \S 3.4）が出発点となる．
\item \textbf{エージェント数と役割設計の拡張}: $n>3$への拡張（Shapley値の近似計算），および「性格」ではなく「推論戦略」（順方向導出・逆方向検証・具体例代入）としてのペルソナ設計の検証．
\item \textbf{ペルソナ注入方式の比較}: 重みSFT・プロンプト・activation steeringの3方式を能力毀損と多様性利得の2軸上で直接比較し，領域に応じた注入方式の選択則を確立すること．
\item \textbf{異種モデル混成と大規模化}: 異なるベースモデルの混成チームや，より大きなモデル規模での2軸トレードオフの定量関係の検証．
\end{itemize}
